\section{Introduction}
\label{sec:intro}

Financial markets are a demanding real-world testbed for autonomous decision-making: errors are not merely incorrect labels, but realized losses after transaction costs, drawdowns, liquidity frictions, and regime shifts. A trading policy must synthesize price dynamics, fundamentals, news, chart structure, and portfolio risk before deciding whether and how much capital to commit. Financial language models established domain understanding and instruction following~\cite{wu2023bloomberggpt,liu2023fingpt}, while surveys trace the transition toward tool-using, memory-augmented agents~\cite{li2023large,zhao2024revolutionizing,wang2024survey}. Recent benchmarks move beyond static question answering to sequential trading, financial-chart comprehension, execution-grounded safety, and reproducible execution assumptions~\cite{stockbench2025,finchartbench2025,finvault2026,yao2026execution}. Together, these advances make finance a consequential setting for agents that must act, learn, and remain auditable under non-stationarity.

Financial agents have progressed along complementary fronts. Role-specialized systems coordinate technical, fundamental, sentiment, and risk perspectives through debate or hierarchical management~\cite{tradingagents2025,fincon2024,hedgeagents2025,finteam2025}. Multimodal and tool-augmented agents connect language reasoning to charts, market data, and executable analytics~\cite{finagent2024}; memory-centric agents preserve observations and reflections~\cite{yu2023finmem,reflexion2023}; and automated research systems generate and screen alpha factors~\cite{rdagent2025,alphaagent2025,factorminer2026,hubble2026}. Recent self-improving designs also feed evaluations into domain knowledge or bounded, human-auditable rule policies~\cite{quantagent2024,sharp2026}. These systems broaden what an agent can observe, organize, and revise, but adaptation remains mediated by external knowledge, prompts or rubrics, retrieved episodes, tools, or factor repositories. The action policy is not directly distilled from dense feedback on its own outcome-matured trajectories, so portfolio-level utility reaches it only indirectly.

Closing this loop is not ordinary supervised fine-tuning. Financial feedback is delayed: no oracle action exists when a trade is placed, and utility becomes observable only after an evaluation horizon closes. The trajectories are also endogenous to the current policy, so a static teacher corpus misses deployment states and errors---the mismatch that motivates on-policy distillation~\cite{gkd2024,seed2026}. Recent OPD variants address unstable or misleading dense transfer with proximal teachers or verifier-based gates~\cite{topd2026,rgopd2026}. In a multi-agent decision, the terminal portfolio outcome is shared even though specialists contribute unequally. Policy parameters, executable factors, and episodic records should also evolve on different clocks: factors must remain versioned and auditable, whereas an episode cannot be retrieved before its outcome is available. Leakage and reproducibility studies reinforce this strict temporal boundary~\cite{profitmirage2025,nguyen2026reliable,yao2026execution}. Together, delayed supervision, endogenous data, shared credit, and asynchronous knowledge states define the learning problem addressed here.

Closed-loop learning becomes well-defined once access to matured outcomes is separated from the time at which each knowledge state may change. We instantiate this separation in \textsc{FinOPD}, a staged framework centered on On-Policy Distillation (OPD). Before OPD, Alpha-Factor Evolution mines, validates, versions, and freezes 157 executable factors in a content-addressed binary artifact. During OPD, a Qwen3.5-9B policy generates the final point-in-time JSON decision from a deterministic specialist trace; only after the horizon closes does a frozen Qwen3.5-27B teacher receive five outcome scalars. Approximate marginal role credit scales final-policy token alignment, the router reuses the top-$k$ masks recorded at rollout, and only accepted high-utility episodes enter availability-filtered memory. The chronology closes the learning loop without exposing the deployed student to future prices, charts, or raw OHLCV paths.
Our contributions are threefold:
\begin{itemize}[leftmargin=*,nosep]
\item We formulate self-evolution in financial multi-agent systems as learning from delayed portfolio outcomes under a strict point-in-time boundary, and identify endogenous current-policy data, shared agent credit, and asynchronously updated knowledge as its defining challenges.
\item We propose \textsc{FinOPD}, which combines a deployable 9B current policy with a training-only 27B hindsight teacher, one explicit cost-aware utility, deterministic coalition credit, stabilized token alignment, a content-addressed 157-factor library, and time-safe memory in one ordered protocol.
\item We provide executable audits for all 157 frozen factors, ordered top-15 routing, time-safe retrieval, the local 9B/27B token update, and a locked factorial evaluator that requires synchronized daily returns, checkpoints, three seeds, block-bootstrap uncertainty, and Holm correction. Missing evidence is surfaced explicitly rather than replaced by a proxy.
\end{itemize}
